\section{Down-conversion}
\label{app:downconversion}

%%%%%%%%%%%%%%%%%%%%%%%%%%%%%%%%%%%%%%%%%%%%%%%%%%%%%%%%%%%%%%%%%%%%%%%%%%%%%%%%%%%%%%%%
%%%%%%%%%%%%%%%%%%%%%%%%%%%%%%%%%%%%%%%%%%%%%%%%%%%%%%%%%%%%%%%%%%%%%%%%%%%%%%%%%%%%%%%%

By the Nyquist Theorem, a signal must be discretely sampled at a rate equal to twice its bandwidth in order to completely represent its information content.
%
Therefore, subject to the finite recording rate of
digital observatory equipment, a radio astronomical signal must be constrained to a limited portion of the radio spectrum.
%
The intermediate process by which the signal from the receiver is band-limited and made ready for discrete sampling is known as down-conversion.

Consider the incoming radio signal, $x(t)$, and its Fourier transform, $X(\nu)$.  The band-limited signal of interest, $x_b(t)$, is parameterized by its centre frequency, $\nu_0$, and bandwidth, \bw.  
%
Baseband down-conversion is the process by which the spectral information originally contained in the range [$\nu_0-\bw/2$, $\nu_0+\bw/2$] is shifted down to [0, \bw]. 

% baseband 1. The original band of frequencies produced by a transducer,
% such as a microphone, telegraph key, or other signal-initiating
% device, prior to initial modulation. Note 1: In transmission systems,
% the baseband signal is usually used to modulate a carrier. Note 2:
% Demodulation re-creates the baseband signal. Note 3: Baseband
% describes the signal state prior to modulation, prior to multiplexing,
% following demultiplexing, and following demodulation. (188) Note 4:
% Baseband frequencies are usually characterized by being much lower in
% frequency than the frequencies that result when the baseband signal is
% used to modulate a carrier or subcarrier. 2. In facsimile, the
% frequency of a signal equal in bandwidth to that between zero
% frequency and maximum keying frequency. (188)

The spectral information is shifted to baseband by demodulating or mixing the radio frequencies (RF) with a local oscillator (LO).  This is equivalent to multiplying the signal, $x(t)$, with a pure tone, $l(t)=a\cos(2\pi\nu_l t + \phi)$.  
%
By application of the convolution theorem, and with reference to \Tab{fourier}, mixing may also be understood as a convolution with a pair of complex-valued delta
functions in the frequency domain.  This understanding proves useful in the following sections.

In addition to mixing to lower frequencies, the signal must also be band-limited before analog-to-digital conversion.  Otherwise, power from frequencies higher than the Nyquist frequency will be reflected back into the band of interest, a pollution known as aliasing.  
%
The ideal low-pass filter is represented by the rectangle function (see \Tab{fourier}) so that a bandpass filter with centre frequency, $\nu_0$, and bandwidth, \bw, is given by
\begin{equation}
\Pi \left( \frac{|\nu|-\nu_0}{\bw} \right).
\end{equation}
Note that the absolute value of $\nu$ in the first term of this equation creates a bandpass window at both positive and negative frequency values.

Down-conversion therefore refers to the combined operation of mixing and band-limiting.  The following two subsections describe in detail two commonly used methods of down-conversion: single-sideband (SSB) and
dual-sideband (DSB).  These are also represented graphically in \Figs{usb}{through}{dsb}.  The process of
down-conversion is performed separately and (ideally) identically on each of the two orthogonal senses of polarization from the receiver feed.

%%%%%%%%%%%%%%%%%%%%%%%%%%%%%%%%%%%%%%%%%%%%
%%%%%%%%%%%%%%%%%%%%%%%%%%%%%%%%%%%%%%%%%%%%
\subsection{Single-sideband Down-conversion}
\label{app:ssb}
%%%%%%%%%%%%%%%%%%%%%%%%%%%%%%%%%%%%%%%%%%%%
%%%%%%%%%%%%%%%%%%%%%%%%%%%%%%%%%%%%%%%%%%%%

During single-sideband down-conversion (SSB), the
signal of interest, $x(t)$, is first bandpass filtered, producing
$x_b(t)$ where \mbox{$X_b(\nu)=X(\nu)\Pi((|\nu|-\nu_0)/\bw)$}.  The
band-limited signal is then mixed with a LO with frequency, $\nu_1$,
producing \mbox{$x_m(t)=x_b(t)\cos(2\pi\nu_1t+\theta)$}, where $\nu_1$ is
set to either 
$\nu_0-\bw/2$ (upper-sideband, shown in \Fig{usb}), producing $x_u(t)$;
or
$\nu_0+\bw/2$ (lower-sideband, shown in \Fig{lsb}), producing $x_l(t)$.
After another stage of low-pass filtering, either $x_u(t)$ or $x_l(t)$
is digitally sampled at the Nyquist rate, $2\bw$.
%
Compared to $x_u(t)$, the two halves of the spectrum in $x_l(t)$
are swapped, which is equivalent to negating the direction on the
frequency axis and taking the complex conjugate of the spectrum.

During playback, the analytic signal associated with $x(t)$ may be
formed in practice by taking the real-to-complex Fast Fourier
Transform (FFT), followed by the complex-to-complex inverse
FFT.  Most real-to-complex FFT implementations automatically omit the
redundant negative frequencies ($X(-\nu)=X^*(\nu)$) from their output,
implicitly producing the analytic signal.
%
Since many signal processing operations
(such as phase-coherent dispersion removal)
are performed in the Fourier domain, the cost of calculating the
analytic signal is transparent.

\begin{figure}
\centerline{\includegraphics[width=0.5\linewidth]{usb.pdf}}
\caption{Upper-sideband down-conversion. From the top, the real-valued signal, $x(t)$
has a complex-valued spectrum with conjugate symmetry, such that $X(-\nu)=X^*(\nu)$,
as depicted by $\real[X(\nu)]$ in blue and $\imag[X(\nu)]$ in red.
The signal is bandpass filtered and the band-limited signal, $X_b(\nu)$, 
is mixed with a local oscillator with frequency
$\nu_0-\bw/2$, as depicted by a pair of blue delta functions.
The resulting signal, $X_m(\nu)$, is low-pass filtered, producing $X_u(\nu)$.
For each spectrum, the frequency axis is drawn in black, and the real and
imaginary parts of each complex value are projected along the blue and red
axes (respectively) depicted at the origin (where $\nu = 0$). }
\label{fig:usb}
\end{figure}

\begin{figure}
\centerline{\includegraphics[width=0.5\linewidth]{lsb.pdf}}
\caption{Lower-sideband down-conversion.  As for upper-sideband
except the band-limited signal, $X_b(\nu)$, is mixed with a local
oscillator with frequency $\nu_0+\bw/2$. The resulting
signal, $X_m(\nu)$, is low-pass filtered, producing $X_l(\nu)$.}
\label{fig:lsb}
\end{figure}


%%%%%%%%%%%%%%%%%%%%%%%%%%%%%%%%%%%%%%%%%%%%
%%%%%%%%%%%%%%%%%%%%%%%%%%%%%%%%%%%%%%%%%%%%
\subsection{Dual-sideband Down-conversion}
\label{app:dsb}

During dual-sideband down-conversion (DSB, see \Fig{dsb}), also known
as quadrature mixing, the voltages from the receiver are split equally
into two signal paths.  One signal is mixed with a local oscillator,
producing
\begin{equation}
	i(t)=x(t)\cos(2\pi\nu_0t).
\end{equation}
The other signal is mixed with the same local oscillator phase-shifted
by 90\deg,
\begin{equation}
q(t)=x(t)\sin(2\pi\nu_0t).
\label{eqn:quadrature}
\end{equation}
Both $i(t)$ and $q(t)$ are low-pass filtered with a cutoff frequency
of $\nu_c=\bw/2$, producing $i_b(t)=i(t)*\sinc(\pi\bw
t)$ and $q_b(t)=q(t)*\sinc(\pi\bw t)$.  The low-pass
filtered signals are then digitally sampled at the Nyquist rate of
$2\nu_c=\bw$.  The signals, $i_b(t)$ and $q_b(t)$ are known as
the in-phase and quadrature components, respectively, of $x(t)$ with
respect to $\nu_0$.

During playback, the analytic signal associated with $x(t)$ is given by
\begin{align}
  z_b(t) & =  i_b(t)+\Ci q_b(t) \nonumber \\
       & =  [x(t)\cos(2\pi\nu_0 t)+\Ci x(t)\sin(2\pi\nu_0 t)] 
		* \sinc(\pi\bw t) \nonumber \\
       & =  [x(t)e^{2\pi\Ci\nu_0t}] * \sinc(\pi\bw t).
\end{align}
%
In the Fourier domain,
\begin{equation}
	Z_b(\nu)=X(\nu+\nu_0) \Pi(\nu/\bw).
\end{equation}
That is, the spectrum of $z_b(t)$ is equivalent to the band-limited
portion of $x(t)$ centred at $\nu_0$.  The negative frequency
components, centred at $-\nu_0$, have been suppressed by low-pass
filtering, forming the analytic signal associated with $x(t)$.

Note that changing the sign of the $90\deg$ phase shift in \eqn{quadrature}
changes the sign of $q(t)$, which results in complex conjugation of $z_b(t)$.
This is equivalent to modulating $x(t)$ with a local oscillator that has
frequency $-\nu_0$, which shifts the negative half of the spectrum of $X(\nu)$
to baseband.  As in the case of lower-sideband down-conversion, 
the resulting spectrum is complex conjugated and the frequency axis is negated.

\begin{figure}
\centerline{\includegraphics[width=\linewidth]{dsb.pdf}}
\caption{Dual-sideband down-conversion (DSB).  The real-valued
signal, $X(\nu)$, is split before mixing with a local oscillator with frequency $\nu_0$ (blue delta functions) and a $90\deg$ phase-shifted local oscillator (red delta functions), 
yielding the in-phase and quadrature components, I$(\nu)$ and Q$(\nu)$, respectively.
%
Each of the signals are low-pass filtered, yielding I$_b(\nu)$ and Q$_b(\nu)$ 
with bandwidth $\bw/2$.  The complex signal,
Z$_b(\nu)$=I$_b(\nu)$+\Ci Q$_b(\nu)$, is the analytic
signal associated with $X_u(\nu)$, which is depicted in \Fig{usb}.}
\label{fig:dsb}
\end{figure}

During SSB down-conversion, bandpass filtering is performed before mixing, necessitating a filter that is tunable over the range of 
frequencies of interest.  In contrast, the low-pass filter used for DSB down-conversion remains constant regardless of the desired observing frequency.  
For this reason, DSB is often the more economical means of down-conversion.

%%%%%%%%%%%%%%%%%%%%%%%%%%%%%%%%%%%%%%%%%%%%
%%%%%%%%%%%%%%%%%%%%%%%%%%%%%%%%%%%%%%%%%%%%
\subsection{Nyquist Zones}
\label{app:Nyquist}
%%%%%%%%%%%%%%%%%%%%%%%%%%%%%%%%%%%%%%%%%%%%
%%%%%%%%%%%%%%%%%%%%%%%%%%%%%%%%%%%%%%%%%%%%

In the examples so far, 
the down-converted signal was placed in the first Nyquist zone, which spans
from $-\bw$ to $\bw$ in the case of single-sideband down-conversion,
and from $-\bw/2$ to $\bw/2$ for dual-sideband down-conversion.
%
The process of discretely sampling the analog signal $x(t)$ at uniformly-spaced instances
in time is equivalent to convolving the band-limited spectrum with an infinite
series of delta functions separated by the sampling frequency.
%
An example of sampling an analog signal after upper-sideband down-conversion
to the first Nyquist zone is depicted in \Fig{first_Nyquist}.

\begin{figure}
\centerline{\includegraphics[width=0.8\linewidth]{nyquist1.pdf}}
\caption{Down-conversion to the first Nyquist zone before sampling.
The band-limited analog signal, $X_{u,1}(\nu)$ is depicted following
upper-sideband down-conversion, along with five of the infinite
series of delta functions separated by the sampling frequency.
The digitized signal $\hat{X}_{u,1}(\nu)$ is the convolution of
the analog signal and the sampling function.}
\label{fig:first_Nyquist}
\end{figure}

It is also possible to place the down-converted and band-limited analog signal in
a region of spectrum offset from the first Nyquist zone.
%
If the spectrum is offset by $N\bw$,
where $N$ is an integer, then the spectrum is said to be placed in the $(N+1)^\mathrm{th}$
Nyquist zone, and the digitized signal will have frequencies that either increase or decrease
linearly across the spectrum.
%
An example of sampling an analog signal after upper-sideband down-conversion
to the second ($N$ = 1) Nyquist zone is depicted in \Fig{second_Nyquist}.
%
With respect to the digitized signal depicted in \Fig{first_Nyquist}, the digitized signal
in \Fig{second_Nyquist} is shifted by $\bw$, which is equivalent to what would
be obtained by digitizing an analog signal after lower-sideband down-conversion
to the first Nyquist zone.
%
That is, the spectrum is complex conjugated and the frequency axis is negated.

\begin{figure}
\centerline{\includegraphics[width=0.8\linewidth]{nyquist2.pdf}}
\caption{Down-conversion to the second Nyquist zone before sampling.
As for \Fig{first_Nyquist}, except that the analog signal has been
down-converted to occupy the second Nyquist zone before analog-to-digital conversion. }
\label{fig:second_Nyquist}
\end{figure}
